% ARCHIVED 2026-10-13 from _c_method.tex: the optional desiderata of P (sparsity, locality, stability) and the
% disabled sigma-consistency definition; unused by any result.

Optionally, the following desiderata may be useful to practitioners:
\begin{itemize}[noitemsep, topsep=-4pt]
  \item \textbf{Desiderata 1} (\textit{Sparsity}) most mass is concentrated on few entries.
  \item \textbf{Desiderata 2} (\textit{Locality}) most per-column mass is concentrated in a short span.
  \item \textbf{Desiderata 3} (\textit{Stability}) small perturbations to $z$ induce small changes in $P$.
\end{itemize}
\iffalse
With provenance, we are ready to define the $\sigma$-consistency constraint from Definition~\ref{def:ProblemSpaceExplanation}:

\definition[\textit{Sigma-Consistency}]
\label{def:SigmaConsistency}
{
Consider some sample $z \in \mathcal{Z}$ in the problem-space, a feature map $\sigma$, provenance $P$, and problem-space explanation $b \in \mathcal{B}$. Then $b$ is said to be consistent under the feature map ($\sigma$-consistent) if and only if $b_i = 0$ for any $i$ in which $P_{i,j} = 0$ for all $j$.
}

This consistency tells us that non-interacting elements should not participate in an explanation. Intuitively, this makes sense and is a useful constraint to place upon $b$.
\fi
